% ============ 第2章 §2.5 Hardy空间的边界值 (I)(原书 pp.134-137) ============
\section{Hardy 空间的边界值 (I)}\label{sec:2.5}
\index{Nevanlinna 类}


如 2.4 节所注意到的,属于 $h^{1}$ 与 $H^{1}$ 的函数利用 Poisson--Stieltjes 积分表示,可以证明几乎处处存在径向极限.属于 $H^{p}$ 的函数,当 $0 < p < 1$ 时有时不能用 Poisson--Stieltjes 积分表示.但本节证明:属于 $H^{p}$ 的函数都几乎处处存在径向极限.

\begin{lemma}\label{lem:2.5.1}
$f \in H(D)$ 且
\begin{equation*}
\sup_{0 \leq r < 1}\int_{-\pi}^{\pi}\log^{+}|f(re^{i\theta})|\dd\theta < \infty,
\end{equation*}
则存在 $h, k \in H^{\infty}$ 使得 $f = h/k$.
\end{lemma}

\begin{proof}
设 $f \in H(D)$ 且 $f(z) \not\equiv 0$,则存在 $m \geq 0$ 使得 $f(z) = \sum_{n=m}^{\infty}a_{n}z^{n}$ 且 $a_{m} \neq 0$.设 $\{z_{n}\}$ 为 $f$ 的非零零点集,且 $0 < |z_{1}| \leq |z_{2}| \leq \cdots < 1$.若 $\{z_{n}\} \cap \{z;\ |z| = \rho < 1\} = \emptyset$,则
\begin{equation*}
F(z) = \log\left\{f(z)\frac{\rho^{m}}{z^{m}}\prod_{|z_{n}| < \rho}\left(\frac{\rho^{2} - \bar{z}_{n}z}{\rho(z - z_{n})}\right)\right\}
\end{equation*}
在 $|z| \leq \rho$ 上正则,且在 $|z| = \rho$ 上 $\operatorname{Re}\left\{F(z)\right\} = \log|f(z)|$.因此由 2.2 节所注意的事实,
\begin{equation*}
F(z) = \frac{1}{2\pi}\int_{-\pi}^{\pi}\log|f(\rho e^{it})|\frac{\rho e^{it} + z}{\rho e^{it} - z}\dd t + iC.
\end{equation*}
因此,当 $\rho < 1$ 时,令
\begin{equation*}
h_{\rho}(z) = \frac{z^{m}}{\rho^{m}}\prod_{|z_{n}| < \rho}\frac{\rho(z - z_{n})}{\rho^{2} - \bar{z}_{n}z}\exp\left\{-\frac{1}{2\pi}\int_{-\pi}^{\pi}\log^{-}|f(\rho e^{it})|\frac{\rho e^{it} + z}{\rho e^{it} - z}\dd t + iC\right\}
\end{equation*}
且
\begin{equation*}
k_{\rho}(z) = \exp\left\{-\frac{1}{2\pi}\int_{-\pi}^{\pi}\log^{+}|f(\rho e^{it})|\frac{\rho e^{it} + z}{\rho e^{it} - z}\dd t\right\},
\end{equation*}
则 $f(z) = h_{\rho}(z)/k_{\rho}(z)$.下面令 $\rho \to 1$ 来证明本引理.

设 $\{\rho_{\ell}\}$ 为满足 $\rho_{\ell} \nearrow 1$ 且 $\{z_{n}\} \cap \{z;\ |z| = \rho_{\ell}\} = \emptyset$ 的数列.令 $H_{\ell}(z) = h_{\rho_{\ell}}(\rho_{\ell}z)$,$K_{\ell}(z) = k_{\rho_{\ell}}(\rho_{\ell}z)$,则 $f(\rho_{\ell}z) = H_{\ell}(z)/K_{\ell}(z)$($z \in D$).$H_{\ell}$ 与 $K_{\ell}$ 在 $D$ 上正则,且 $|H_{\ell}(z)| \leq 1$,$|K_{\ell}(z)| \leq 1$($z \in D$).因此 $\{H_{\ell}\}$ 与 $\{K_{\ell}\}$ 是正规族,故对任意含于 $D$ 的闭集 $E$,存在数列 $\{\ell_{j}\}$ 使 $H_{\ell_{j}}$ 与 $K_{\ell_{j}}$ 在 $E$ 上一致收敛.设其极限为 $h$ 与 $k$,由一致定理,$h$ 与 $k$ 与 $E$ 无关,在 $D$ 上正则,$|h(z)| \leq 1$ 且 $|k(z)| \leq 1$.因此令 $j \to \infty$,应当有 $f(z) = h(z)/k(z)$($z \in D$),但必须证明 $k(z) \not\equiv 0$.这里才首次用到关于 $f$ 的条件
\begin{equation*}
\sup_{0 \leq r < 1}\int_{-\pi}^{\pi}\log^{+}|f(re^{i\theta})|\dd\theta < \infty.
\end{equation*}
由于 $\log^{+}|f(\rho_{\ell}z)| \geq -\log|K_{\ell}(z)|$,利用 Jensen 公式(命题~\ref{prop:2.2.1})与 $f$ 的条件,得 $|K_{\ell}(0)| \geq \varepsilon > 0$,所以 $k(z) \not\equiv 0$.
\end{proof}

\begin{lemma}\label{lem:2.5.2}
若 $k \in H^{\infty}$ 且 $k(z) \not\equiv 0$,则 $k$ 几乎处处存在径向极限 $k(e^{i\theta})$,且 $\log|k(e^{i\theta})| \in L^{1}$.
\end{lemma}

\begin{proof}
可以设 $\|k\|_{\infty} \leq 1$.由 $k(z) \not\equiv 0$,可以写 $k(z) = z^{m}k_{0}(z)$,$m \geq 0$ 且 $k_{0}(0) \neq 0$,故可以设 $k(0) \neq 0$.由 Jensen 公式(命题~\ref{prop:2.2.1}),
\begin{equation*}
\frac{1}{2\pi}\int_{-\pi}^{\pi}\log|k(re^{i\theta})|\dd\theta = \log|k(0)| + \sum_{|z_{n}| < r}\log\frac{r}{|z_{n}|}
\end{equation*}
成立.这里 $\{z_{n}\}$ 是 $k$ 的零点集.因此
\begin{equation*}
\frac{1}{2\pi}\int_{-\pi}^{\pi}\log|k(re^{i\theta})|\dd\theta
\end{equation*}
关于 $r$ 单调递增.由于 $k$ 属于 $H^{1}$,如 2.4 节所注意到的,几乎处处存在径向极限 $k(e^{i\theta})$,故由 Lebesgue 积分的 Fatou 引理,
\begin{equation*}
\int_{-\pi}^{\pi}\left|\log|k(e^{i\theta})|\right|\dd\theta \leq \varliminf_{r \to 1}\left\{-\int_{-\pi}^{\pi}\log|k(re^{i\theta})|\dd\theta\right\}.
\end{equation*}
上式右边的积分单调递减,故为有限值,$\log|k(e^{i\theta})| \in L^{1}$.
\end{proof}

\begin{theorem}\label{thm:2.5.3}
当 $0 < p \leq \infty$ 时,若 $f \in H^{p}$ 且 $f(z) \not\equiv 0$,则 $f$ 几乎处处存在径向极限 $f(e^{i\theta})$.此时 $\log|f(e^{i\theta})| \in L^{1}$ 且 $f(e^{i\theta}) \in L^{p}$.
\end{theorem}

\begin{proof}
当 $0 < p < \infty$ 时,由 $\lim_{x \to \infty}(\log x)/x^{p} = 0$,若 $f \in H^{p}$ 则
\begin{equation*}
\sup_{0 \leq r < 1}\int_{-\pi}^{\pi}\log^{+}|f(re^{i\theta})|\dd\theta < \infty.
\end{equation*}
因此由引理~\ref{lem:2.5.1},若 $f(z) \not\equiv 0$,则存在 $h, k \in H^{\infty}$ 使得 $f = h/k$.由引理~\ref{lem:2.5.2},$h$ 与 $k$ 几乎处处存在径向极限 $h(e^{i\theta})$ 与 $k(e^{i\theta})$,且 $k(e^{i\theta})$ 在正测度集上不为零,所以径向极限 $f(e^{i\theta})$ 存在.由于 $f \in H^{p}$,由 Fatou 引理知 $f(e^{i\theta}) \in L^{p}$.又由引理~\ref{lem:2.5.2},$\log|h(e^{i\theta})| \in L^{1}$ 且 $\log|k(e^{i\theta})| \in L^{1}$,故 $\log|f(e^{i\theta})| \in L^{1}$.
\end{proof}

引理~\ref{lem:2.5.1} 的逆也成立.用 Jensen 公式证明逆命题是容易的.满足引理~\ref{lem:2.5.1} 的条件的 $H(D)$ 的函数的全体称为 Nevanlinna 类 $N$,且 $N \supset \bigcup_{p > 0}H^{p}$.
